% Vietnamese translation for OI Public Library
% Source-File: IOI中国国家候选队论文/2005/杨思雨/杨思雨.doc
\documentclass[11pt]{article}
\usepackage{oipl-vn}

\newcommand{\Oh}{\mathcal{O}}
\newcommand{\floor}[1]{\left\lfloor #1\right\rfloor}

\title{Vẻ đẹp ở khắp nơi: bàn về liên hệ giữa tỉ lệ vàng và tin học}
\author{Yang Siyu}
\date{2005}

\begin{document}
\maketitle

\origtitle{Vẻ đẹp ở khắp nơi: bàn về liên hệ giữa tỉ lệ vàng và tin học}
\sourcepath{IOI China candidate team papers / 2005 / Yang Siyu.doc}

\renewcommand{\abstractname}{Tóm tắt}
\begin{abstract}
Bài viết giới thiệu mối liên hệ giữa tỉ lệ vàng và các kì thi lập trình, cũng
như ứng dụng của nó trong một số bài toán tin học. Nội dung được chia thành
bốn phần: phần mở đầu trình bày nguồn gốc của tỉ lệ vàng và sự xuất hiện của
nó trong nhiều lĩnh vực; phần thứ hai phân tích dần một ví dụ để chỉ ra bản
chất của bài toán có liên hệ với số tỉ lệ vàng; phần thứ ba giới thiệu phương
pháp tối ưu hóa theo tỉ lệ vàng qua một ví dụ khác; phần cuối tổng kết rằng
tỉ lệ vàng có quan hệ khá mật thiết với tin học thi đấu.
\end{abstract}

\noindent\textbf{Từ khóa:} tỉ lệ vàng; tin học; liên hệ.

\section{Mở đầu}

Ngay từ thời Hy Lạp cổ đại, con người đã phát hiện ra tỉ lệ vàng. Hơn hai
nghìn năm trước, nhà toán học Eudoxus nêu ra bài toán sau: trên đoạn thẳng
\(AB\) có điểm \(P\), chia \(AB\) thành hai đoạn \(AP\) và \(PB\). Nếu tỉ số
độ dài của hai đoạn này đúng bằng tỉ số giữa cả đoạn \(AB\) và đoạn dài hơn
\(AP\), thì điểm \(P\) phải nằm ở vị trí nào trên \(AB\)?

Giả sử độ dài \(AP\) bằng \(1\), độ dài \(PB\) bằng \(x\). Khi đó độ dài
\(AB\) là \(1+x\). Từ điều kiện
\[
  \frac{1}{x}=\frac{1+x}{1},
\]
ta được phương trình
\[
  x^2+x-1=0.
\]
Nghiệm dương của phương trình là
\[
  x=\frac{\sqrt{5}-1}{2},
\]
kí hiệu là \(\varphi\), có giá trị xấp xỉ \(0.618\). Tương ứng, độ dài đoạn
\(AB\) là
\[
  1+x=\frac{\sqrt{5}+1}{2},
\]
kí hiệu là \(\Phi\).

Về sau, con người nghiên cứu rất nhiều quanh tỉ số này. Nghệ sĩ và nhà khoa
học Leonardo da Vinci còn đặt tên cho \(\varphi\) là ``số tỉ lệ vàng''. Điều
thú vị là tỉ lệ vàng xuất hiện khắp nơi trong tự nhiên và trong nhiều lĩnh
vực khác. Chẳng hạn, rốn người nằm gần điểm chia vàng của toàn bộ chiều cao
cơ thể; đầu gối lại nằm gần điểm chia vàng của đoạn từ rốn đến gót chân. Ở
một số loài thực vật, góc giữa hai cuống lá liền kề là \(137^\circ 28'\), gần
đúng bằng góc tạo bởi hai bán kính chia đường tròn theo tỉ lệ \(1:0.618\).

Ngoài ra, tỉ lệ vàng cũng thường được xem là biểu tượng của vẻ đẹp. Người ta
cho rằng các đối tượng có tỉ lệ phù hợp với tỉ lệ vàng thường hài hòa và giàu
tính thẩm mĩ. Vì vậy, tỉ lệ này có nhiều ứng dụng trong kiến trúc, hội họa,
điêu khắc, nhiếp ảnh và âm nhạc. Ví dụ, trong thiết kế đền Parthenon thời Hy
Lạp cổ đại có nhiều nhóm tỉ lệ vàng; trong bức tranh thuyền cá bên bờ của
Vincent van Gogh, chiều rộng và chiều cao toàn bức tranh lần lượt được cột
buồm và đường chân trời chia thành hai phần theo tỉ lệ gần với tỉ lệ vàng.

Thật ra, trong tin học cũng ẩn chứa những điều tương tự. Bài viết này thông
qua hai ví dụ để giới thiệu một số liên hệ giữa tin học và tỉ lệ vàng.

\section{Ví dụ 1: Trò chơi lấy đá}

\subsection{Phát biểu}

Có hai đống đá. Sau khi trò chơi bắt đầu, hai người luân phiên lấy đá. Mỗi
lượt có hai cách lấy: lấy một số viên tùy ý từ một trong hai đống, hoặc lấy
cùng một số viên từ cả hai đống. Người lấy hết những viên đá cuối cùng là
người thắng. Cho trước số viên ban đầu của hai đống là \(a\) và \(b\), với
\(0\le a,b\le 10^9\). Giả sử cả hai bên đều chơi tối ưu, hãy xác định người
đi trước có chiến lược thắng hay không.

\subsection{Thuật toán 1}

Trong bài toán, hai số \(a\) và \(b\) quyết định thắng thua cuối cùng. Do đó
ta dùng \((a,b)\) để biểu diễn trạng thái hiện tại. Bằng cách xây dựng cây
trò chơi, ta có thể phán định trạng thái ban đầu \((a,b)\) là trạng thái
thắng hay thua. Cần chú ý rằng hai đống đá không phân biệt vai trò, nên
\((a,b)\) và \((b,a)\) thực chất là cùng một trạng thái.

Xét ví dụ ban đầu có hai đống lần lượt là \(1\) và \(2\). Khi xây dựng cây
trò chơi từ trạng thái \((1,2)\), các trạng thái con trùng lặp chỉ cần mở rộng
một lần. Sau đó có thể phân loại thắng thua theo ba trường hợp:
\begin{enumerate}
  \item Nếu một trạng thái không có trạng thái con, dùng trực tiếp luật chơi
  để phán định thắng thua.
  \item Nếu một trạng thái có ít nhất một trạng thái con là thua cho người
  đi trước, thì trạng thái đó là thắng cho người đi trước.
  \item Nếu mọi trạng thái con đều là thắng cho người đi trước, thì trạng
  thái đó là thua cho người đi trước.
\end{enumerate}

Nhờ vậy ta suy ra được trạng thái ban đầu thắng hay thua. Khi cài đặt, có thể
dùng quy hoạch động hoặc tìm kiếm có ghi nhớ. Thuật toán này có độ phức tạp
không gian \(\Oh(N^2)\) và độ phức tạp thời gian \(\Oh(N^3)\).

\subsection{Thuật toán 2}

Thuật toán trên rõ ràng chưa giải quyết bài toán một cách hiệu quả. Ta cần
khai thác sâu hơn đặc điểm của luật chơi. Từ hai cách lấy đá, có thể thấy nếu
trạng thái \((a,b)\) là thua cho người đi trước, thì các trạng thái sau đều
phải là thắng cho người đi trước:
\begin{itemize}
  \item \((a,i)\), với \(i\ne b\), hoặc \((i,b)\), với \(i\ne a\);
  \item \((a+i,b+i)\), với \(i\ne 0\).
\end{itemize}

Nói cách khác, mỗi số nguyên dương xuất hiện đúng một lần trong toàn bộ các
trạng thái thua; đồng thời hiệu giữa hai đống đá của hai trạng thái thua bất
kì cũng khác nhau. Dựa trên đặc điểm này, ta có thể xây dựng tất cả trạng
thái thua.

Ban đầu, trạng thái thua nhỏ nhất là \((0,0)\). Theo quy luật trên, trong
những trạng thái thua tiếp theo sẽ không còn trạng thái có một đống bằng
\(0\), và cũng không có trạng thái mà hai đống bằng nhau. Vì vậy, ở trạng thái
thua thứ hai, đống nhỏ hơn phải có số đá là số nguyên dương nhỏ nhất chưa
xuất hiện, tức \(1\); hiệu giữa hai đống cũng phải là hiệu nhỏ nhất chưa xuất
hiện, tức \(1\). Do đó trạng thái thua thứ hai là \((1,2)\).

Từ đây ta có cách xây dựng mọi trạng thái thua: lúc đầu chỉ có \((0,0)\). Ở
lần xây dựng thứ \(i\), tìm số nguyên dương nhỏ nhất \(a\) chưa xuất hiện
trong các trạng thái thua đã biết; trạng thái thua mới là \((a,a+i)\). Thuật
toán này có độ phức tạp thời gian và không gian đều là \(\Oh(N)\).

\subsection{Thuật toán 3}

Dù thuật toán 2 đã cải thiện rất nhiều so với thuật toán 1, kích thước dữ
liệu của bài toán vẫn quá lớn. Khi việc phân tích rơi vào bế tắc, ta thử
quan sát các trường hợp nhỏ.

\begin{center}
\begin{tabular}{c|c|c|c}
\textbf{Thứ tự \(i\)} & \textbf{Trạng thái thua \((a,b)\)} &
\textbf{\(a/i\)} & \textbf{\(b/i\)} \\
\hline
1 & \((1,2)\) & 1.0000 & 2.0000 \\
2 & \((3,5)\) & 1.5000 & 2.5000 \\
3 & \((4,7)\) & 1.3333 & 2.3330 \\
4 & \((6,10)\) & 1.5000 & 2.5000 \\
5 & \((8,13)\) & 1.6000 & 2.6000 \\
\(\cdots\) & \(\cdots\) & \(\cdots\) & \(\cdots\) \\
997 & \((1613,2610)\) & 1.6179 & 2.6179 \\
998 & \((1614,2612)\) & 1.6172 & 2.6172 \\
999 & \((1616,2615)\) & 1.6176 & 2.6176 \\
1000 & \((1618,2618)\) & 1.6180 & 2.6180 \\
\end{tabular}
\end{center}

Trong bảng trên, có thể thấy với các trạng thái được xây dựng về sau, tỉ số
giữa số đá \(a\) của đống nhỏ hơn và số thứ tự \(i\) rất gần \(\Phi\). Điều
này gợi ý rằng \(a\) có thể có liên hệ với \(\Phi i\). Vì vậy ta quan sát
ngược lại quan hệ giữa \(\Phi i\) và \(a\).

\begin{center}
\begin{tabular}{c|c|c}
\textbf{Thứ tự \(i\)} & \textbf{\(\Phi i\)} &
\textbf{Trạng thái \((a,b)\)} \\
\hline
1 & 1.6180 & \((1,2)\) \\
2 & 3.2361 & \((3,5)\) \\
3 & 4.8541 & \((4,7)\) \\
4 & 6.4721 & \((6,10)\) \\
5 & 8.0902 & \((8,13)\) \\
\(\cdots\) & \(\cdots\) & \(\cdots\) \\
997 & 1613.1799 & \((1613,2610)\) \\
998 & 1614.7979 & \((1614,2612)\) \\
999 & 1616.4160 & \((1616,2615)\) \\
1000 & 1618.0340 & \((1618,2618)\) \\
\end{tabular}
\end{center}

Trong các trường hợp nhỏ ở bảng thứ hai, với trạng thái thua thứ \(i\), số
đá của đống nhỏ hơn đều đúng bằng phần nguyên của \(\Phi i\). Ta đưa ra phỏng
đoán:
\[
  \text{trạng thái thua thứ } i
  =
  \bigl(\floor{\Phi i},\floor{(\Phi+1)i}\bigr).
\]
Phỏng đoán này có thể chứng minh là đúng; chứng minh được trình bày trong
phụ lục. Do đó, với trạng thái \((a,b)\), giả sử \(a\le b\) và đặt
\(i=b-a\). Nếu
\[
  a=\floor{\Phi i}
  \quad\text{và}\quad
  b=\floor{(\Phi+1)i},
\]
thì người đi trước chắc chắn thua; ngược lại người đi trước chắc chắn thắng.

Như vậy, độ phức tạp thời gian và không gian đều là hằng số, và bài toán được
giải quyết trọn vẹn.

\subsection{Nhận xét}

Bảng sau so sánh ba thuật toán đã nêu.

\begin{center}
\begin{tabular}{c|c|c|c}
\textbf{Thuật toán} & \textbf{Thời gian} & \textbf{Không gian} &
\textbf{Tư tưởng chính} \\
\hline
1 & \(\Oh(N^3)\) & \(\Oh(N^2)\) & Quy hoạch động trên cây trò chơi \\
2 & \(\Oh(N)\) & \(\Oh(N)\) & Xây dựng các trạng thái thua \\
3 & \(\Oh(1)\) & \(\Oh(1)\) & Phán định trực tiếp thắng thua \\
\end{tabular}
\end{center}

Thuật toán 1 là phương pháp thường gặp khi xử lí bài toán trò chơi, đồng thời
giúp định nghĩa rõ trạng thái. Thuật toán 2 khai thác sâu hơn luật chơi để
tìm ra tính chất của các trạng thái thua. Thuật toán 3 tiếp tục phân tích dữ
liệu nhỏ và tìm được thuật toán hằng số. Qua quá trình phân tích và đào sâu
như vậy, ta thấy bản chất của bài toán có liên hệ với tỉ lệ vàng, và phát
hiện vẻ đẹp ẩn trong một bài toán trò chơi.

\section{Ví dụ 2: Bài toán leo núi}

\subsection{Phát biểu}

Cho một hàm đơn đỉnh \(f(x)\) có giá trị cực đại trên đoạn \([0,L]\). Cần tìm
hoành độ \(x_0\) tại đó hàm đạt cực đại bằng một chuỗi truy vấn. Mỗi lần
truy vấn, chương trình đưa một điểm \(t\) cho thư viện tương tác, và thư viện
trả về \(f(t)\).

Giới hạn dữ liệu là \(0<L\le 10^4\). Số lần truy vấn không được vượt quá
\(40\), và kết quả tính được phải sai khác điểm tối ưu \(x_0\) không quá
\(10^{-3}\).

\subsection{Phân tích}

Với hai điểm truy vấn khác nhau trong đoạn, gọi điểm có giá trị hàm lớn hơn
là điểm tốt, và điểm có giá trị hàm nhỏ hơn là điểm xấu. Ta sẽ chỉ ra rằng
điểm tối ưu và điểm tốt luôn nằm cùng phía so với điểm xấu.

Trước hết, điểm tối ưu hiển nhiên không trùng với điểm xấu. Nếu điểm tối ưu
trùng với điểm tốt thì kết luận là hiển nhiên. Nếu điểm tối ưu không nằm ở
điểm tốt, có hai khả năng. Trường hợp thứ nhất, điểm tốt và điểm xấu nằm ở
hai phía của điểm tối ưu; khi đó kết luận vẫn hiển nhiên. Trường hợp thứ hai,
cả hai điểm truy vấn nằm cùng phía so với điểm tối ưu; vì \(f(x)\) là hàm đơn
đỉnh, điểm truy vấn gần điểm tối ưu hơn chắc chắn là điểm tốt, nên điểm tối
ưu và điểm tốt vẫn nằm cùng phía so với điểm xấu.

Kết luận trên gợi ý rằng ta có thể so sánh giá trị hàm tại hai điểm trong
đoạn mục tiêu, rồi loại bỏ một phần đoạn mục tiêu. Cứ như vậy, phạm vi tìm
kiếm được thu nhỏ dần và tiến sát điểm tối ưu cho đến khi đạt sai số cho
phép.

Trong bài gốc, sai số cho phép bằng một phần mười triệu độ dài đoạn ban đầu.
Để bảo đảm độ chính xác, ta cần thu nhỏ đoạn mục tiêu còn \(1/10^8\) so với
ban đầu. Vấn đề cần nghiên cứu là làm thế nào thu nhỏ phạm vi nhanh nhất.
Điểm then chốt nằm ở cách chọn vị trí truy vấn.

Giả sử đoạn mục tiêu hiện tại là \([a,b]\), hai điểm truy vấn là \(x_1\) và
\(x_2\), với \(x_1<x_2\). Trước khi truy vấn, ta không biết trước điểm nào sẽ
là điểm tốt hay điểm xấu; xác suất để mỗi điểm trở thành điểm xấu là như
nhau. Vì vậy, khi bố trí điểm truy vấn, hai điểm nên đối xứng qua trung điểm
của đoạn hiện tại:
\[
  x_1-a=b-x_2.
\]
Đây là nguyên tắc thứ nhất, gọi là nguyên tắc đối xứng.

Nguyên tắc đối xứng rất thường gặp. Ví dụ, tìm kiếm nhị phân mà ta hay dùng
trong các bài toán tìm kiếm cũng hoàn toàn tuân theo nguyên tắc này. Từ đó ta
dễ nghĩ đến cách dùng nhị phân cho bài toán: mỗi lần chọn hai điểm rất gần
nhau quanh trung điểm đoạn hiện tại, so sánh giá trị hàm ở hai điểm đó, rồi
xác định phạm vi chứa điểm tối ưu. Cách này mỗi lần có thể loại bỏ gần một
nửa độ dài đoạn hiện tại. Tuy nhiên, để bảo đảm độ chính xác, số lần truy vấn
cần thiết xấp xỉ \(2\log_2 10^8\), khoảng \(54\) lần, vượt quá yêu cầu của
đề.

Nguyên nhân số truy vấn quá lớn là mỗi lần muốn loại bỏ một nửa đoạn, ta phải
thực hiện lại hai truy vấn. Trong khi đó, điểm tốt cũ vẫn nằm trong đoạn hiện
tại, nhưng nhị phân không tận dụng thông tin giá trị hàm tại điểm ấy. Vì vậy,
ta nghĩ đến việc mỗi vòng chỉ truy vấn thêm một điểm trong đoạn hiện tại, rồi
so sánh nó với điểm tốt cũ để thu được điểm tốt và điểm xấu mới, từ đó tiếp
tục thu nhỏ phạm vi.

Do mỗi lần chọn điểm truy vấn đều phải thỏa mãn nguyên tắc đối xứng, vị trí
hai điểm truy vấn ban đầu sẽ quyết định độ dài phần đoạn có thể bị loại bỏ ở
các vòng sau. Nếu đặt \(x_1\) và \(x_2\) rất gần nhau như trong nhị phân, lần
đầu có thể loại bỏ gần \(50\%\) đoạn \([a,b]\), nhưng những lần sau chỉ loại
bỏ được một phần rất nhỏ. Điều đó không có lợi cho việc nhanh chóng tiến sát
điểm tối ưu. Ta rút ra nguyên tắc thứ hai: tốt nhất là mỗi lần loại bỏ cùng
một tỉ lệ của đoạn hiện tại. Có thể gọi đây là nguyên tắc loại bỏ theo tỉ lệ.

Theo hai nguyên tắc trên, xét lần truy vấn thứ nhất và thứ hai tại \(x_1\) và
\(x_2\), với \(x_1<x_2\). Sau lần so sánh đầu tiên, do tính đối xứng, độ dài
phần bị loại bỏ trong hai khả năng đều là \(b-x_2\). Không mất tính tổng quát,
giả sử \(x_1\) là điểm tốt, \(x_2\) là điểm xấu, nên đoạn bị loại bỏ là
\((x_2,b)\). Lần truy vấn thứ ba đặt tại \(x_3\). Khi đó \(x_3\) cần nằm trong
\([a,x_2]\) và đối xứng với \(x_1\) trong đoạn này. Đồng thời \(x_3\) phải
nằm bên trái \(x_1\); nếu không, sau khi so sánh \(x_3\) với \(x_1\), phần bị
loại bỏ sẽ có cùng độ dài như lần trước chứ không giữ cùng tỉ lệ, trái với
nguyên tắc loại bỏ theo tỉ lệ.

Vì vậy, sau khi so sánh \(x_3\) với \(x_1\), độ dài phần bị loại bỏ đều bằng
\(x_2-x_1\). Theo nguyên tắc loại bỏ theo tỉ lệ, ta có
\[
  \frac{b-x_2}{b-a}
  =
  \frac{x_2-x_1}{x_2-a},
  \qquad
  x_1-a=b-x_2.
\]
Biến đổi và rút gọn hai đẳng thức trên cho ta
\[
  \frac{x_2-a}{b-a}
  =
  \frac{\sqrt{5}-1}{2}
  =\varphi.
\]
Đẳng thức này rất giống phương trình chia đoạn ở phần mở đầu. Nó cho thấy
\(x_2\) nằm tại điểm chia vàng của đoạn \([a,b]\). Nói cách khác, đặt điểm
truy vấn tại điểm chia vàng của đoạn là lựa chọn hợp lí. Phương pháp xấp xỉ
tối ưu này vì thế được gọi là phương pháp tỉ lệ vàng.

\subsection{Hiệu quả}

Vì mỗi điểm truy vấn đều nằm tại điểm chia vàng của đoạn mục tiêu hiện tại,
tỉ lệ giữa đoạn được giữ lại và đoạn cũ trong trường hợp xấu nhất là
\(\varphi:1\). Do đó số lần truy vấn xấp xỉ
\[
  2+\log_{1/\varphi} 10^8,
\]
vừa đủ nằm trong giới hạn \(40\) lần, nên bài toán được giải quyết hoàn chỉnh.

Trong bài toán này, ta thoát khỏi khuôn mẫu của nhị phân, tận dụng đầy đủ
thông tin của mỗi lần truy vấn và đưa ra hai nguyên tắc khi chọn điểm truy
vấn. Chính tính chất tỉ lệ đặc biệt của tỉ lệ vàng phù hợp với nguyên tắc
loại bỏ theo tỉ lệ, từ đó mở ra lời giải cho bài toán.

\section{Tổng kết}

Hai ví dụ trên đều có lời giải liên hệ mật thiết với số tỉ lệ vàng.

Trong ví dụ 1, các cặp số được xây dựng ẩn chứa quy luật của tỉ lệ vàng. Các
kì thi lập trình thường xuất hiện những bài toán có màu sắc toán học mạnh; mà
số tỉ lệ vàng lại có quan hệ gần gũi với việc xây dựng dãy số, ma trận và các
cấu trúc tương tự, nên nó thường trở thành chìa khóa cho loại bài này.

Trong ví dụ 2, tính chất tỉ lệ đặc biệt của tỉ lệ vàng trở thành điểm then
chốt của lời giải. Trong các bài toán truy vấn tương tự, tỉ lệ vàng có ứng
dụng khá rộng; về mặt lí thuyết có thể chứng minh rằng phương pháp tỉ lệ vàng
vừa giới thiệu dùng số lần truy vấn tối ưu trong các phương pháp tìm kiếm một
chiều kiểu này.

Dĩ nhiên, liên hệ giữa tỉ lệ vàng và tin học thi đấu không chỉ dừng ở đây;
vẫn còn nhiều điều đáng để khám phá. Nhà điêu khắc Auguste Rodin từng nói:
vẻ đẹp ở khắp nơi; đối với đôi mắt của chúng ta, điều thiếu không phải là vẻ
đẹp, mà là sự phát hiện. Thật vậy, chỉ cần chăm suy nghĩ và dám khám phá, mỗi
người đều có thể có một đôi mắt biết phát hiện vẻ đẹp.

\section{Tài liệu tham khảo}

\begin{enumerate}
  \item Jiang Lu chủ biên, \textit{Bách khoa toàn thư thiếu niên nhi đồng
  Trung Quốc: khoa học và kĩ thuật}, tái bản lần 2, Nhà xuất bản Giáo dục
  Chiết Giang, 1994.12.
  \item Tobias Dantzig, Su Zhongxiang dịch, \textit{Số: ngôn ngữ của khoa
  học}, bản thứ nhất, Nhà xuất bản Giáo dục Thượng Hải, 2000.12.
  \item Nhóm biên soạn \textit{Sổ tay toán học}, \textit{Sổ tay toán học},
  bản thứ nhất, Nhà xuất bản Giáo dục Nhân dân, 1979.05.
  \item Zhang Weiyong chủ biên, \textit{Đề thi và lời giải Olympic Tin học
  thanh thiếu niên, tỉnh An Huy 1994--2004}, bản thứ nhất, Nhà xuất bản Đại
  học Công nghiệp Hợp Phì, 2004.08.
  \item Luo Ji, \textit{Phân tích hai hướng tư duy khi giải bài toán đối sách:
  từ bài lấy đá}, luận văn đội tuyển quốc gia năm 2002, 2002.02.
  \item Hua Luogeng, \textit{Tối ưu học}, bản thứ nhất, Nhà xuất bản Khoa học,
  1981.04.
\end{enumerate}

\appendix

\section{Chứng minh phỏng đoán trong ví dụ 1}

Trong phần này, \(\floor{x}\) biểu thị phần nguyên của \(x\), còn \(\{x\}\)
biểu thị phần thập phân của \(x\).

\paragraph{Định nghĩa 1.1.}
Theo thuật toán 2, xây dựng ma trận \(A\) có hai cột sao cho các trạng thái
thua lần lượt là
\[
  (A_{1,1},A_{1,2}),\ (A_{2,1},A_{2,2}),\
  (A_{3,1},A_{3,2}),\ldots
\]

\paragraph{Định nghĩa 1.2.}
Theo phỏng đoán, xây dựng ma trận \(B\) có hai cột:
\[
  B_{i,1}=\floor{\Phi i},
  \qquad
  B_{i,2}=\floor{(\Phi+1)i}.
\]

\paragraph{Bổ đề 1.1, định lí Beatty.}
Nếu \(a\) và \(b\) là hai số vô tỉ dương thỏa mãn
\[
  \frac{1}{a}+\frac{1}{b}=1,
\]
và
\[
  P=\{x\mid x=\floor{at},\ t\in\mathbb{Z}^{+}\},\qquad
  Q=\{x\mid x=\floor{bt},\ t\in\mathbb{Z}^{+}\},
\]
thì \(P\) và \(Q\) tạo thành một phân hoạch của tập số nguyên dương
\(\mathbb{Z}^{+}\).

\paragraph{Chứng minh.}
Với số nguyên dương bất kì \(n\), giả sử trong tập \(P\) có \(x\) phần tử nhỏ
hơn \(n\), và trong tập \(Q\) có \(y\) phần tử nhỏ hơn \(n\). Khi đó
\[
  x=\floor{\frac{n}{a}},\qquad
  y=\floor{\frac{n}{b}},
\]
vì \(a\) và \(b\) vô tỉ nên \(n/a\) và \(n/b\) không phải số nguyên. Do đó
\[
  x+y
  =
  \floor{\frac{n}{a}}+\floor{\frac{n}{b}}
  <
  \frac{n}{a}+\frac{n}{b}
  =
  n,
\]
đồng thời
\[
  n-(x+y)
  =
  \left\{\frac{n}{a}\right\}+\left\{\frac{n}{b}\right\}
  <2.
\]
Suy ra \(n-2<x+y<n\). Vì \(x+y\) là số nguyên, \(x+y=n-1\). Như vậy, với mọi
\(n\), trong hai tập \(P,Q\) có đúng \(n-1\) phần tử nhỏ hơn \(n\), và có đúng
\(n\) phần tử nhỏ hơn \(n+1\). Do đó phần tử bằng \(n\) xuất hiện trong đúng
một trong hai tập. Vậy mỗi số nguyên dương thuộc đúng một trong \(P\) và
\(Q\), tức \(P\cup Q=\mathbb{Z}^{+}\) và \(P\cap Q=\varnothing\).

\paragraph{Hệ quả 1.1.}
Nếu \(a\) và \(b\) là hai số vô tỉ dương thỏa mãn
\[
  \frac{1}{a}+\frac{1}{b}=1,
\]
thì mỗi số nguyên dương \(x\) có đúng một biểu diễn dưới dạng \(\floor{at}\)
hoặc \(\floor{bt}\), với \(t\in\mathbb{Z}^{+}\).

\paragraph{Hệ quả 1.2.}
Mỗi số nguyên dương xuất hiện đúng một lần trong ma trận \(B\).

\paragraph{Chứng minh.}
Ta có
\[
  \frac{1}{\Phi}+\frac{1}{\Phi+1}=1.
\]
Theo hệ quả 1.1 và định nghĩa của \(B\), kết luận được chứng minh.

\paragraph{Bổ đề 1.2.}
Với mọi \(i>1\), \(B_{i,1}\) là số nguyên dương nhỏ nhất chưa xuất hiện trong
\(i-1\) hàng đầu của ma trận \(B\).

\paragraph{Chứng minh.}
Theo định nghĩa 1.2, mọi số nguyên nhỏ hơn \(B_{i,1}=\floor{\Phi i}\) đều phải
xuất hiện trong các hàng trước hàng \(i\). Cụ thể, các số ở cột thứ nhất có
dạng \(\floor{\Phi t}\); số phần tử nhỏ hơn \(\floor{\Phi i}\) đến từ cột
này là số giá trị \(t<i\). Tương tự, các số ở cột thứ hai có dạng
\(\floor{(\Phi+1)t}\). Dùng hệ quả 1.2, các số nhỏ hơn \(\floor{\Phi i}\)
không thể bị trùng giữa hai cột, nên toàn bộ các số nguyên dương nhỏ hơn
\(\floor{\Phi i}\) đều đã nằm trong \(i-1\) hàng đầu. Vì bản thân
\(\floor{\Phi i}\) nằm ở hàng \(i\), nó chưa xuất hiện trước đó. Do đó
\(B_{i,1}\) đúng là số nguyên dương nhỏ nhất chưa xuất hiện trong các hàng
trước.

\paragraph{Định lí 1.}
Hai ma trận \(A\) và \(B\) tương đương. Với mọi \(i\),
\[
  A_{i,1}=B_{i,1},\qquad A_{i,2}=B_{i,2}.
\]

\paragraph{Chứng minh.}
Ta quy nạp theo số hàng \(i\). Khi \(i=1\), hiển nhiên
\[
  A_{1,1}=B_{1,1},\qquad A_{1,2}=B_{1,2}.
\]
Giả sử với mọi \(i<k\) đều có \(A_{i,1}=B_{i,1}\) và \(A_{i,2}=B_{i,2}\).
Theo định nghĩa 1.1, \(A_{k,1}\) là số nguyên dương nhỏ nhất chưa xuất hiện
trong \(k-1\) hàng đầu của \(A\). Theo bổ đề 1.2, \(B_{k,1}\) cũng là số
nguyên dương nhỏ nhất chưa xuất hiện trong \(k-1\) hàng đầu của \(B\). Từ giả
thiết quy nạp suy ra hai tập số đã xuất hiện trước hàng \(k\) là như nhau, vì
thế
\[
  A_{k,1}=B_{k,1}.
\]
Theo cách xây dựng của thuật toán 2,
\[
  A_{k,2}=A_{k,1}+k.
\]
Mặt khác, từ định nghĩa 1.2,
\[
  B_{k,2}
  =
  \floor{(\Phi+1)k}
  =
  \floor{\Phi k}+k
  =
  B_{k,1}+k.
\]
Do đó \(A_{k,2}=B_{k,2}\). Định lí được chứng minh, và phỏng đoán về các
trạng thái thua trong thuật toán 3 cũng đúng.

\section{Định nghĩa hàm đơn đỉnh trong ví dụ 2}

Theo đề bài, ở đây chỉ cần định nghĩa hàm đơn đỉnh \(f(x)\) có cực đại. Giả
sử \(x_0\) là điểm đạt giá trị lớn nhất của hàm trên đoạn \([a,b]\). Khi đó:
\begin{itemize}
  \item với mọi \(x_1<x_2<x_0\), ta có
  \(f(x_1)<f(x_2)<f(x_0)\);
  \item với mọi \(x_0<x_1<x_2\), ta có
  \(f(x_0)>f(x_1)>f(x_2)\).
\end{itemize}

\section{Thư viện tương tác trong ví dụ 2}

Thí sinh C++ dùng thư viện \texttt{tools\_c.h}. Thí sinh Pascal dùng thư viện
\texttt{tools\_p.ppu}. Bản gốc gửi lời cảm ơn Zhou Dong từ Trường Trung học
số 1 thành phố Wuhu, tỉnh An Huy, vì đã viết thư viện tương tác cho bài này.

\section{Mô tả chương trình}

Phụ lục của bản gốc kèm bốn chương trình Pascal. Vì các chương trình này chỉ
minh họa trực tiếp các thuật toán đã phân tích ở trên và trong mã nguồn có
nhiều chú thích tiếng Trung, bản dịch nêu ngắn vai trò của từng chương trình
thay vì chép nguyên văn.

\begin{description}
  \item[\texttt{Stone\_1.pas}] Cài đặt thuật toán 1 bằng tìm kiếm đệ quy có
  ghi nhớ trên các trạng thái \((a,b)\), đồng thời tận dụng tính đối xứng
  giữa \((a,b)\) và \((b,a)\).
  \item[\texttt{Stone\_2.pas}] Cài đặt thuật toán 2, đánh dấu các số nguyên
  đã xuất hiện và xây dựng dần các trạng thái thua.
  \item[\texttt{Stone\_3.pas}] Cài đặt thuật toán 3, dùng hằng
  \(\Phi=(\sqrt{5}+1)/2\) để phán định trực tiếp trạng thái thua.
  \item[\texttt{Explorer.pas}] Cài đặt phương pháp tỉ lệ vàng cho bài toán
  leo núi, mỗi vòng giữ lại điểm truy vấn tốt cũ và hỏi thêm một điểm đối
  xứng trong đoạn hiện tại.
\end{description}

\end{document}
